\subsection{诱导过程与转移结构：显式有限标记图}\label{subsec:Phi_m-graph}

下面给出 $Y_m$ 的一个完全显式的 sofic 图表示，从而刻画“将 $\Fold_m$ 沿长轨迹连续重叠窗口施加后，输出过程的状态/转移结构”。

\begin{definition}[长度 \texorpdfstring{$m-1$}{m-1} 后缀状态图]\label{def:G_m}
令顶点集为
\[
V_m:=\{0,1\}^{m-1}.
\]
对任意顶点 $v=v_1\cdots v_{m-1}\in V_m$ 与任意 $b\in\{0,1\}$，定义一条有向边
\[
e=(v,b):\ v\to v':=v_2\cdots v_{m-1}b,
\]
并用 $\ell_m(e)\in X_m$ 标记该边：
\[
\ell_m(e):=\Fold_m(v_1\cdots v_{m-1}b)\in X_m.
\]
由此得到一个有限标记图 $G_m=(V_m,E_m,\ell_m)$，其中 $|V_m|=2^{m-1}$，$|E_m|=2^{m}$。
\leanverified{phiState\_card}
\end{definition}

\begin{theorem}[显式 sofic 表示与可计算诱导结构]\label{thm:Phi_m-sofic-graph}
令 $Z_m\subset E_m^{\ZZ}$ 为图 $G_m$ 的双边边移位（即双边路径空间），并令
\[
(\mathcal{L}_m(e))_t:=\ell_m(e_t)\in X_m
\]
为逐边取标记得到的序列。则
\[
Y_m=\mathcal{L}_m(Z_m),
\]
从而 $Y_m$ 是一个由 $G_m$ 给出的 sofic 子移位；并且 $G_m$ 给出一个可审计的有限状态诱导结构：状态为长度 $m-1$ 的输入后缀，转移为追加一比特，输出为对应窗口经 $\Fold_m$ 的稳定类型。
\leanverified{paper\_phi\_m\_sofic\_graph}
\leanverified{phiEdge\_card}
\leanverified{phiGraph\_outdegree\_two}
\end{theorem}

\begin{proof}
\textbf{（包含 $\subset$）} 给定任意输入序列 $s\in\{0,1\}^{\ZZ}$，定义
\[
v_t:=s_t s_{t+1}\cdots s_{t+m-2}\in V_m,\qquad
e_t:=(v_t,s_{t+m-1})\in E_m.
\]
则 $e=(e_t)_{t\in\ZZ}$ 是 $G_m$ 上的一条双边路径（因为 $v_{t+1}$ 恰为 $v_t$ 去首加尾的更新）。并且由标记定义，
\[
(\mathcal{L}_m(e))_t=\ell_m(e_t)=\Fold_m(s_t\cdots s_{t+m-1})=(\Phi_m(s))_t.
\]
因此 $\Phi_m(s)\in \mathcal{L}_m(Z_m)$，即 $Y_m\subset \mathcal{L}_m(Z_m)$。

\textbf{（包含 $\supset$）} 反之，取任意双边路径 $e=(e_t)\in Z_m$。每条边 $e_t$ 具有形式 $(v_t,b_t)$，且路径条件保证 $v_{t+1}$ 等于 $v_t$ 去首加尾的更新。由此可一致地构造一个二元序列 $s\in\{0,1\}^{\ZZ}$：令 $s_{t+m-1}:=b_t$，并用 $v_t$ 提供窗口的其余 $m-1$ 位（路径相容性保证该定义不矛盾）。于是对该 $s$ 有
\[
(\Phi_m(s))_t=\Fold_m(s_t\cdots s_{t+m-1})=\Fold_m(v_t b_t)=\ell_m(e_t)=(\mathcal{L}_m(e))_t.
\]
因此 $\mathcal{L}_m(e)\in Y_m$，从而 $\mathcal{L}_m(Z_m)\subset Y_m$。两边相等，证毕。
\end{proof}

\begin{corollary}[序列层折叠的熵不降：每步保留 \texorpdfstring{$1$}{1} 比特自由度]\label{cor:Phi_m-entropy-no-drop}
对任意 \(m\ge 2\)，令 \(Z_m\subset E_m^{\ZZ}\) 为图 \(G_m\) 的双边边移位、\(\mathcal{L}_m:Z_m\to X_m^{\ZZ}\) 为逐边取标记映射（定义 \ref{def:G_m}），并记 \(Y_m=\mathcal{L}_m(Z_m)\)（定理 \ref{thm:Phi_m-sofic-graph}）。
则：
\begin{enumerate}
  \item \(Z_m\) 与二元全移位 \(\{0,1\}^{\ZZ}\) 共轭，因而 \(h_{\mathrm{top}}(Z_m)=\log 2\)；
  \item 标记图 \(G_m\) 右解析（同一顶点出发的两条出边标记不同；定理 \ref{thm:pom-right-resolving}），从而 \(\mathcal{L}_m\) 是有限对一因子码（每个输出序列的 lift 数至多为 \(|V_m|=2^{m-1}\)）；
  \item 因而拓扑熵保持：
  \[
  \boxed{\ h_{\mathrm{top}}(Y_m)=h_{\mathrm{top}}(Z_m)=\log 2\ }.
  \]
\end{enumerate}
该结论给出“窗口层多对一压缩”与“序列层记忆补偿”的严格命题：\(\Phi_m\) 将窗口层丢失的纤维自由度迁移到时间方向的有限状态记忆中，而不在熵意义下压缩每步 \(1\) 比特的自由度。
\leanverified{LabeledPath}
\leanverified{labeledPath\_initial\_injective}
\leanverified{labeledPath\_card\_le\_state\_card}
\leanverified{labeledPath\_count\_le\_states}
\leanverified{paper\_labeled\_path\_lift\_bound}
\leanverified{paper\_Phi\_m\_entropy\_no\_drop}
\end{corollary}
\begin{proof}
\textbf{(1)}--\textbf{(2)} 见定理 \ref{thm:Phi_m-sofic-graph}（共轭）与右解析性；\textbf{(3)} 有限对一因子码熵保持可参见 \cite{LindMarcus1995,Kitchens1998SymbolicDynamics}。
\end{proof}

\begin{remark}[时间记忆作为隐寄存器]\label{rem:Phi-m-memory-implicit-register}
窗口层的多对一压缩并不等价于时间层的信息损失：当 $m\ge 3$ 时，定理 \ref{thm:Phi_m-conjugacy-threshold} 给出有限窗逆码 $\Psi_m$，从而被折叠抹去的自由度可迁移到时间方向的有限状态记忆并在有限延迟后恢复。这提供了“第三寄存器”的另一种实现方式：其不必以静态轴显式存储，而可作为同步滤波器的内部状态隐式承载。在规范群胚口径下，该隐寄存器对应局部截面之间的粘接数据（见小节 \ref{subsec:pom-gauge-groupoid}）。
\end{remark}


\subsubsection{可逆性阈值：\texorpdfstring{$\Phi_m$}{Phi\_m} 的共轭与纤维坍缩}\label{subsec:Phi_m-conjugacy}

上一节已将 $Y_m$ 组织为标记图 $G_m$ 的像（定理 \ref{thm:Phi_m-sofic-graph}），并证明熵率不降（推论 \ref{cor:Phi_m-entropy-no-drop}）。
本节给出可逆性与纤维结构的确定结论：在 Zeckendorf 折叠 $\Fold_m$ 的口径下，除去 $m=2$ 的两点退化外，$\Phi_m$ 在 $m\ge 3$ 时为共轭。

\paragraph{可终止的局部判定口径（子集滤波）}
沿定义 \ref{def:G_m}，对任意候选后缀集合 $S\subseteq V_m$ 与输出符号 $a\in X_m$，定义一步更新
$$
\delta_m(S,a):=\{v'\in V_m:\exists\,e:v\to v'\ \text{s.t. }v\in S,\ \ell_m(e)=a\}.
$$
从初始集合 $S_0:=V_m$ 出发，对任意输出块 $a_0a_1\cdots a_{t-1}$ 递推
$$
S_{j+1}:=\delta_m(S_j,a_j)\qquad(0\le j<t).
$$
该递推可视为“仅由输出块驱动”的有限状态滤波器：$S_j$ 记录读入前 $j$ 个输出符号后，内部后缀（顶点）仍可能取到的候选集。
当 $S_t$ 为单点集时，意味着该长度-$t$ 输出块已将内部后缀同步到唯一值。

\begin{theorem}[可逆性阈值与有限窗逆码]\label{thm:Phi_m-conjugacy-threshold}
令 $\Phi_m$ 为定义 \ref{def:Phi_m} 的滑动块码，$Y_m$ 为定义 \ref{def:Y_m} 的像子移位。
则：
\begin{enumerate}
  \item 当 $m\ge 3$ 时，$\Phi_m:\{0,1\}^{\ZZ}\to Y_m$ 为拓扑共轭；更具体地，存在一个滑动块码
  $$
  \Psi_m:Y_m\to\{0,1\}^{\ZZ},
  $$
  其\textbf{记忆}为 $m-1$、\textbf{预期性}为 $0$，并满足
  $$
  \Psi_m\circ\Phi_m=\mathrm{id}_{\{0,1\}^{\ZZ}},\qquad
  \Phi_m\circ\Psi_m=\mathrm{id}_{Y_m}.
  $$
  \item 当 $m=2$ 时，$\Phi_2$ 恰好在两条常值序列 $0^{\ZZ}$ 与 $1^{\ZZ}$ 上出现 $2$-对-$1$ 退化，其余点一一。
\end{enumerate}
\end{theorem}
\leanverified{paper\_Phi\_m\_conjugacy\_threshold}

\begin{corollary}[像子移位的共轭类塌缩（\(m\ge 3\)）]\label{cor:Ym-conjugacy-all-mge3}
对任意 $m,n\ge 3$，移位系统 $(Y_m,\sigma)$ 与 $(Y_n,\sigma)$ 拓扑共轭。
更具体地，若取定理 \ref{thm:Phi_m-conjugacy-threshold} 中的逆码 $\Psi_m:Y_m\to\{0,1\}^{\ZZ}$，则
\[
h_{m\to n}:=\Phi_n\circ \Psi_m:\ Y_m\longrightarrow Y_n
\]
为同胚且满足 \(h_{m\to n}\circ\sigma=\sigma\circ h_{m\to n}\)。
\end{corollary}
\begin{proof}
由定理 \ref{thm:Phi_m-conjugacy-threshold}，当 \(m\ge 3\) 时 \(\Phi_m\) 为 \(\{0,1\}^{\ZZ}\) 与 \(Y_m\) 的共轭，其逆映射为 \(\Psi_m\)。
因此 \(h_{m\to n}\) 为同胚；又由于 \(\Phi_n\) 与 \(\Psi_m\) 均与移位交换，故 \(h_{m\to n}\) 与移位交换。
证毕。
\end{proof}
\leanverified{paper\_Ym\_conjugacy\_all\_mge3}
\leanverified{paper\_ym\_conjugacy\_all\_mge3}

\begin{proof}
证明思路建立在“有限可判定”的滤波器上。

\paragraph{局部同步：长度-$m$ 块把候选集压到单点}
对固定 $m$，考虑上述滤波更新 $S_{j+1}=\delta_m(S_j,a_j)$，并记
$$
R_m:=\{S_m:\exists\,a_0\cdots a_{m-1}\in\mathcal{L}_m(Y_m)\ \text{s.t. 递推得到 }S_m\}.
$$
若 $R_m$ 中每个元素都是单点集，则任意长度-$m$ 的输出块都把内部后缀同步到唯一顶点。
该命题对每个固定 $m$ 是一个\textbf{完全有限}的判定问题：状态空间为 $V_m$ 的幂集可达部分，转移由 $G_m$ 给出。
脚本 \texttt{scripts/exp\_phi\_m\_conjugacy\_threshold.py} 对 $m\in[2,12]$ 的范围给出可复算证书，输出汇总见表 \ref{tab:phi_m_conjugacy_threshold}。
\begin{table}[H]
\centering
\scriptsize
\setlength{\tabcolsep}{6pt}
\caption{Conjugacy threshold certificate for the sliding-block factor $\Phi_m$. We build the subset-filter automaton from the labeled graph $G_m$ and compute $R_{m-1}$ and $R_m$, the sets of reachable filter states after exactly $m-1$ and $m$ output symbols. If all states in $R_m$ are singletons, then each length-$m$ output word has a unique length-$(2m-1)$ preimage block, yielding a finite-window inverse (memory $m-1$). The non-singleton count in $R_{m-1}$ provides a sharpness witness for the sync delay. We also cross-check periodic-point counts up to a small $n$.}
\label{tab:phi_m_conjugacy_threshold}
\begin{tabular}{r r r r r r r l r}
\toprule
$m$ & det states & det edges & $|R_{m-1}|$ & non-singleton ends $(m-1)$ & $|R_m|$ & non-singleton ends $(m)$ & periodic $\#\mathrm{Fix}_n$ & checked $n$\\
\midrule
2 & 3 & 7 & 3 & 1 & 3 & 1 & $2^n-1$ & 8\\
3 & 7 & 20 & 5 & 1 & 4 & 0 & $2^n$ & 8\\
4 & 16 & 50 & 10 & 2 & 8 & 0 & $2^n$ & 8\\
5 & 32 & 95 & 20 & 4 & 16 & 0 & $2^n$ & 8\\
6 & 69 & 204 & 40 & 8 & 32 & 0 & $2^n$ & 8\\
7 & 142 & 404 & 80 & 16 & 64 & 0 & $2^n$ & 8\\
8 & 296 & 827 & 160 & 32 & 128 & 0 & $2^n$ & 8\\
9 & 607 & 1656 & 320 & 64 & 256 & 0 & $2^n$ & 8\\
10 & 1244 & 3342 & 640 & 128 & 512 & 0 & $2^n$ & 8\\
11 & 2532 & 6703 & 1280 & 256 & 1024 & 0 & $2^n$ & 8\\
12 & 5141 & 13464 & 2560 & 512 & 2048 & 0 & $2^n$ & 8\\
\bottomrule
\end{tabular}
\end{table}

其中 $m\ge 3$ 时 $R_m$ 的“non-singleton ends $(m)$”为 $0$，而 $m=2$ 时在任何步数都出现唯一一个非单点端态，对应常值输出块所导致的两点退化。
同时表中 “non-singleton ends $(m-1)$” 在 $m\ge 3$ 时严格为正，这给出同步延迟的尖锐性见证：一般需要完整的长度-$m$ 输出块才能把候选集压到单点（对应逆码记忆为 $m-1$）。

\paragraph{从单点端态到唯一 lift：利用双向解析性}
由定理 \ref{thm:pom-right-resolving} 与定理 \ref{thm:pom-left-resolving}，标记图 $G_m$ 在 Zeckendorf 折叠下同时右解析与左解析：同一顶点出发（或到达）的一步扩展可由标记唯一判别。
因此一旦长度-$m$ 输出块把端态同步为单点 $\{v'\}$，则该块在 $G_m$ 上的长度-$m$ lift 路径不仅存在，而且由左解析性可从终点 $v'$ 逐步唯一回溯得到；从而该输出块对应唯一的长度-$(2m-1)$ 输入块。

\paragraph{构造逆码 \texorpdfstring{$\Psi_m$}{Psi\_m}}
当 $m\ge 3$ 时，由上述两步可知：对任意允许的长度-$m$ 输出块
$$
a_{t-m+1}\cdots a_t\in\mathcal{L}_m(Y_m),
$$
存在唯一的 $G_m$-路径
$$
e_{t-m+1},\dots,e_t\in E_m
$$
使得其逐边标记恰为该输出块。
令 $e_{t-m+1}=(v,b)$（其中 $b\in\{0,1\}$ 是该边对应的追加比特），定义
$$
(\Psi_m(y))_t:=b,
$$
其中 $y\in Y_m$ 且 $a_j=y_j$。
该定义仅依赖 $y_{t-m+1},\dots,y_t$，故 $\Psi_m$ 的记忆为 $m-1$、预期性为 $0$。
由构造立得 $\Psi_m\circ\Phi_m=\mathrm{id}$，并因 $\Phi_m$ 满射到 $Y_m$ 得 $\Phi_m\circ\Psi_m=\mathrm{id}_{Y_m}$。

\paragraph{$m=2$ 的两点退化}
脚本证书给出：唯一的非单点端态对应常值输出块，其恰来自两条输入常值序列 $0^{\ZZ},1^{\ZZ}$；
此外对任意其它输出点滤波在有限步内同步到单点，从而 lift 唯一。
综上得到 $m=2$ 的“仅两点 $2$-对-$1$”退化结论。
\end{proof}

\begin{remark}[可逆性阈值的同调解释：\(m=2\) 的唯一分支点对应不可消去的 \texorpdfstring{$\mathbb{Z}/2$}{Z/2} 隐变量]\label{rem:Phi_m-conjugacy-threshold-z2-homology}
定理 \ref{thm:Phi_m-conjugacy-threshold} 可被理解为：一个 \(\mathbb{Z}/2\)-值的纤维标签是否能被长度-\(m\) 的输出窗口制度性读出，从而被“规范化”为同调零（coboundary）的判定。
滤波器递推 \(S_{j+1}=\delta_m(S_j,a_j)\) 记录在给定输出块下内部后缀的候选集合；当 \(m\ge 3\) 时，任意长度-\(m\) 输出块将候选集合同步到单点（表 \ref{tab:phi_m_conjugacy_threshold}），因而该 \(\mathbb{Z}/2\) 标签在窗口内可被完全同步并由逆码 \(\Psi_m\) 消去，表现为共轭。
相反，当 \(m=2\) 时，唯一的非单点端态恰对应两条常值轨道 \(0^{\ZZ},1^{\ZZ}\)，等价于：在该分支轨道上存在一个不可由长度-\(2\) 输出决定的 \(\mathbb{Z}/2\) 隐变量；它给出唯一的拓扑分支点并构成共轭失败的唯一障碍。
\end{remark}

\begin{corollary}[\texorpdfstring{$\zeta_{Y_m}$}{zetaYm} 的封口与“折叠因子不产出非平凡 prime shadow”]\label{cor:Phi_m-zeta-closed}
对 $m\ge 3$，由于 $Y_m$ 与二元全移位共轭（定理 \ref{thm:Phi_m-conjugacy-threshold}），其周期点计数满足
$$
\#\mathrm{Fix}(\sigma^n|_{Y_m})=2^n\qquad(n\ge 1),
$$
从而内禀 Artin--Mazur \(\zeta\)-函数为
$$
\zeta_{Y_m}(z)=\frac{1}{1-2z}.
$$
而 $m=2$ 时有
$$
\#\mathrm{Fix}(\sigma^n|_{Y_2})=2^n-1,\qquad
\zeta_{Y_2}(z)=\frac{1-z}{1-2z}.
$$
因此当 $m\ge 3$ 时，$Y_m$ 自身的 prime-orbit/\(\zeta\) 指纹退化为全移位骨架；本文中出现的非平凡 prime shadow 必然来自”核 + 势/权重 + 扭曲通道”（而非单纯来自折叠因子 $Y_m$）。
\leanverified{paper\_fold\_image\_subshift\_zeta\_closed}
\end{corollary}


\begin{remark}[关于 Markov 性与可计算近似]\label{rem:Phi_m-not-markov}
定理 \ref{thm:Phi_m-sofic-graph} 给出的是\emph{拓扑/组合}层面的有限状态表示：$Y_m$ 的所有允许序列由有限图刻画。但对任意输入过程（例如第 \ref{sec:experiments} 节的旋转扫描产生的 Sturmian 序列），其推前测度 $(\Phi_m)_*\mu$ 在 $Y_m$ 上未必是一阶马尔可夫测度。若需要“诱导转移结构”的数值近似，可在图 $G_m$ 上用长样本估计边/状态频率，从而得到一个可审计的有限状态 Markov 近似；该近似与第 \ref{sec:quant-fit} 节的直方图/偏差指标对齐。
\end{remark}
